\chapter{Application to Arrhythmia Detection}

% use Times
%\usepackage{times}
% For figures
%\usepackage{graphicx} % more modern

% For Tables
%\usepackage{booktabs}      % professional-quality tables

% For citations
%\usepackage{natbib}

% For algorithms
%\usepackage{algorithm}
%\usepackage{algorithmic}
%\usepackage[caption = false]{subfig}
%\usepackage{hyperref}
%\usepackage{array}

% Packages hyperref and algorithmic misbehave sometimes.  We can fix
% this with the following command.
%\newcommand{\theHalgorithm}{\arabic{algorithm}}

% TODO figure out labelling
\section{Introduction}
\label{sec:arrhythmias:intro}

We develop a model which can diagnose irregular heart rhythms, also known as
arrhythmias, from single-lead ECG signals better than a cardiologist. Key to
exceeding expert performance is a deep convolutional network which can map a
sequence of ECG samples to a sequence of arrhythmia annotations along with a
novel dataset two orders of magnitude larger than previous datasets of its
kind.

Many heart diseases, including Myocardial Infarction, AV Block, Ventricular
Tachycardia and Atrial Fibrillation can all be diagnosed from ECG signals with
an estimated 300 million ECGs recorded annually \cite{heden1996detection}. We
investigate the task of arrhythmia detection from the ECG record. This is known
to be a challenging task for computers but can usually be determined by an
expert from a single, well-placed lead.

Arrhythmia detection from ECG recordings is usually performed by expert
technicians and cardiologists given the high error rates of computerized
interpretation.  One study found that of all the computer predictions for
non-sinus rhythms, only about 50\% were correct \cite{shah2007errors}; in
another study, only 1 out of every 7 presentations of second degree AV block
were correctly recognized by the algorithm \cite{guglin2006common}. To
automatically detect heart arrhythmias in an ECG, an algorithm must implicitly
recognize the distinct wave types and discern the complex relationships between
them over time. This is difficult due to the variability in wave morphology
between patients as well as the presence of noise.

We train a 34-layer convolutional neural network (CNN) to detect arrhythmias in
arbitrary length ECG time-series. Figure~\ref{fig:record} shows an example of
an input to the model. In addition to classifying noise and the sinus rhythm,
the network learns to classify and segment twelve arrhythmia types present in
the time-series. The model is trained end-to-end on a single-lead ECG signal
sampled at 200Hz and a sequence of annotations for every second of the ECG as
supervision. To make the optimization of such a deep model tractable, we use
residual connections and batch-normalization \cite{he2016deep, ioffe2015batch}.
The depth increases both the non-linearity of the computation as well as the
size of the context window for each classification decision.

We construct a dataset 500 times larger than other datasets of its kind
\cite{moody2001impact, goldberger2000physiobank}. One of the most popular
previous datasets, the MIT-BIH corpus contains ECG recordings from 47 unique
patients. In contrast, we collect and annotate a dataset of about 30,000 unique
patients from a pool of nearly 300,000 patients who have used the Zio Patch
monitor\footnote[1]{iRhythm Technologies, San Francisco, California}
\cite{turakhia2013diagnostic}. We intentionally select patients exhibiting
abnormal rhythms in order to make the class balance of the dataset more even
and thus the likelihood of observing unusual heart-activity high.

We test our model against board-certified cardiologists. A committee of three
cardiologists serve as gold-standard annotators for the 336 examples in the
test set. Our model exceeds the individual expert performance on both recall
(sensitivity), and precision (positive predictive value) on this test set.

\section{Related Work}
\label{sec:arrhythmias:related}

Automatic high-accuracy methods for R-peak extraction have existed at least
since the mid 1980's~\cite{pan1985real}. Current algorithms for R-peak
extraction tend to use wavelet transformations to compute features from the raw
ECG followed by finely-tuned threshold based classifiers~\cite{li1995detection,
martinez2004wavelet}. Because accurate estimates of heart rate and heart rate
variability can be extracted from R-peak features, feature-engineered
algorithms are often used for coarse-grained heart rhythm classification,
including detecting tachycardias (fast heart rate), bradycardias (slow heart
rate), and irregular rhythms. However, such features alone are not sufficient
to distinguish between most heart arrhythmias since features based on the
atrial activity of the heart as well as other features pertaining to the QRS
morphology are needed.

Much work has been done to automate the extraction of other features from the
ECG. For example, beat classification is a common sub-problem of
heart-arrhythmia classification. Drawing inspiration from automatic speech
recognition, Hidden Markov models with Gaussian observation probability
distributions have been applied to the task of beat
detection~\cite{coast1990approach}. Artificial neural networks have also been
used for the task of beat detection~\cite{melo2000arrhythmia}. While these
models have achieved high-accuracy for some beat types, they are not yet
sufficient for high-accuracy heart arrhythmia classification and segmentation.
For example,~\cite{artis1991detection} train a neural network to distinguish
between Atrial Fibrillation and Sinus Rhythm on the MIT-BIH dataset. While the
network can distinguish between these two classes with high-accuracy, it does
not generalize to noisier single-lead recordings or classify among the full
range of 15 rhythms available in MIT-BIH. This is in part due to insufficient
training data, and because the model also discards critical information in the
feature extraction stage.

The most common dataset used to design and evaluate ECG algorithms is the
MIT-BIH arrhythmia database~\cite{moody2001impact} which consists of 48
half-hour strips of ECG data. Other commonly used datasets include the MIT-BIH
Atrial Fibrillation dataset~\cite{moody1983new} and the QT
dataset~\cite{laguna1997database}. While useful benchmarks for R-peak
extraction and beat-level annotations, these datasets are too small for
fine-grained arrhythmia classification. The number of unique patients is in the
single digit hundreds or fewer for these benchmarks. A recently released
dataset captured from the AliveCor ECG monitor contains about 7000
records~\cite{clifford2017}.  These records only have annotations for Atrial
Fibrillation; all other arrhythmias are grouped into a single bucket. The
dataset we develop contains 29,163 unique patients and 14 classes with hundreds
of unique examples for the rarest arrhythmias.

Machine learning models based on deep neural networks have consistently been
able to approach and often exceed human agreement rates when large annotated
datasets are available~\cite{amodei2016deep, xiong2016achieving,he2015delving}.
These approaches have also proven to be effective in healthcare applications,
particularly in medical imaging where pretrained ImageNet models can be
applied~\cite{esteva2017dermatologist, gulshan2016development}. We draw on work
in automatic speech recognition for processing time-series with deep
convolutional neural networks and recurrent neural
networks~\cite{hannun2014deepspeech, sainath2013deep}, and techniques in deep
learning to make the optimization of these models tractable~\cite{he2016deep,
he2016identity, ioffe2015batch}.


%\begin{figure}[ht]
%  \centering
%  \includegraphics[width=\linewidth]{figs/main_figure.png}
%  \caption{
%    Our trained convolutional neural network correctly detecting the sinus rhythm (SINUS) and Atrial Fibrillation (AFIB) from this ECG recorded with a single-lead wearable heart monitor.
%  }
%  \label{fig:record}
%\end{figure}
%\chapter{Introduction}

\section{Overview}
\section{Contributions}
\section{Outline}
\section{First Published Appearances}

%
%\section{Model}
%\label{model}
%\begin{figure}[ht!]
%  \centering
%  \includegraphics[height=15cm]{figs/ecg-nnet.pdf}
%  \caption{The architecture of the network. The first and last layer are special-cased due to the pre-activation residual blocks. Overall, the network contains 33 layers of convolution followed by a fully-connected layer and a softmax.}
%  \label{fig:net}
%\end{figure}
%
%
%\section{Model}
\label{sec:first_pass:model}

We train neural networks using the CTC loss function to do maximum likelihood
training of letter sequences given acoustic features as input. We consider a
single utterance as a training example consisting of an acoustic feature matrix
$X$ and word transcription $Y$. The CTC objective function maximizes the log
probability $\log p(Y \mid X)$.

\subsection{Deep Neural Networks}
\label{sec:first_pass:model:dnn}

With the loss function fixed we must next define how we compute $p(c \mid
x_t)$, the predicted distribution over output characters $c$ given the audio
features $x_t$ at time $t$. While many function approximators are possible for
this task, we choose as our most basic model a DNN. A DNN computes the
distribution $p(c \mid x_t)$ using a series of hidden layers followed by an
output layer. Given an input vector $x_t$ the first hidden layer activations
are a vector computed as,
\begin{align*}
  h^{(1)} = \sigma(W^{(1)} x_t + b^{(1)}).
\end{align*}
The matrix $W^{(1)}$ and vector $b^{(1)}$ are the weight matrix and bias vector
for the layer. The function $\sigma(\cdot)$ is a point-wise nonlinearity. We
use rectifier nonlinearities and thus choose, $\sigma(z) = \max (z, 0)$.

DNNs can have arbitrarily many hidden layers. After the first hidden
layer, the hidden activations $h^{(l)}$ for layer $l$ are computed as,
\begin{align*}
  h^{(l)} = \sigma(W^{(l)} h^{(l-1)} + b^{(l)}).
\end{align*}

To obtain a proper distribution over the set of possible characters, the
final layer of the network is a \emph{softmax} output layer of the form,
\begin{align*}
  p(c=c_k \mid x_t) = \frac{\exp(- ( W^{(s)}_k h^{(L)} + b^{(s)}_k))}
  {\sum_j \exp(- ( W^{(s)}_j h^{(L)} + b^{(s)}_j))},
\end{align*}
where $W^{(s)}_k$ is the $k$'th row of the output weight matrix $W^{(s)}$,
$b^{(s)}_k$ is a scalar bias term and $L$ is the index of the last hidden
layer.

We can compute a subgradient for all parameters of the DNN given a training
example and thus utilize gradient-based optimization techniques. Note that this
same DNN formulation is commonly used in DNN-HMM models to predict a
distribution over senones instead of characters.

\subsection{Recurrent Deep Neural Networks}

A transcription $W$ has many temporal dependencies which a DNN may not
sufficiently capture. At each timestep $t$ the DNN computes its output using
only the input features $x_t$, ignoring previous hidden representations and
output distributions. To enable better modeling of the temporal dependencies
present in a problem, we use a RNN. We select one hidden layer $l$
to have a temporally recurrent weight matrix $W^{(f)}$ and compute the layer's
hidden activations as,
\begin{align*}
  h^{(l)}_t = \sigma(W^{(l)} h^{(l-1)}_t +  W^{(f)} h^{(l)}_{t-1} + b^{(l)}).
\end{align*}
Note that we now make the distinction $h^{(l)}_t$ for the hidden activation
vector of layer $l$ at timestep $t$ since it now depends upon the activation
vector of layer $l$ at time $t-1$.

When working with RNNs, we found it important to use a modified version of the
rectifier nonlinearity. This modified function selects $\sigma(z) = \min( \max
(z, 0), 20)$ which clips large activations to prevent divergence during network
training. Setting the maximum allowed activation to 20 results in the clipped
rectifier acting as a normal rectifier function in all but the most extreme
cases.

Aside from these changes, computations for a RNN are the same as those in a DNN
as described in \ref{sec:first_pass:model:dnn}. Like the DNN, we can compute a
subgradient for a RNN using backpropagation through time. In our experiments we
always compute the gradient completely through time rather than truncating to
obtain an approximate subgradient.

\subsection{Bi-Directional Recurrent Deep Neural Networks}

While forward recurrent connections reflect the temporal nature of the audio
input, a more powerful sequence transduction model is a bidirectional RNN (BRNN), which
maintains state both forwards and backwards in time. Such a model can integrate
information from the entire temporal extent of the input features when making
each prediction. We extend the RNN to form a BRNN by again choosing a
temporally recurrent layer $l$. The BRNN creates both a forward and backward
intermediate hidden representation which we call $h_t^{(f)}$ and $h_t^{(b)}$
respectively. 

We use the temporal weight matrices $W^{(f)}$ and $W^{(b)}$ to propagate
$h_t^{(f)}$ forward in time and $h_t^{(b)}$ backward in time respectively. We
update the forward and backward components via the equations,
\begin{align*}
    h^{(f)}_t &= \sigma(W^{(j)T} h^{(j-1)}_t +  W^{(f)T} h^{(f)}_{t-1} + b^{(j)}), \\
    h^{(b)}_t &= \sigma(W^{(j)T} h^{(j-1)}_t +  W^{(b)T} h^{(b)}_{t+1} + b^{(j)}).
\end{align*}

Note that the recurrent forward and backward hidden representations are
computed entirely independently from each other. As with the RNN we use the
modified nonlinearity function $\sigma(z) = \min( \max (z,0), 20)$. To obtain
the final representation $h^{(l)}_t$ for the layer we sum the two temporally
recurrent components,
\begin{align*}
  h^{(l)}_t = h^{(f)}_t + h^{(b)}_t.
\end{align*}
Aside from this change to the recurrent layer the BRNN computes its output
using the same equations as the RNN. As for other models, we can compute a
subgradient for the BRNN directly to perform gradient-based optimization.

%
%\begin{figure}
%  \centering
%  \includegraphics[width=\linewidth]{figs/results}
%  \caption{
%    Evaluated on the test set, the model outperforms the average cardiologist score on both the Sequence and the Set F1 metrics.
%  }
%  \label{fig:model_vs_cardiol}
%\end{figure}
%
%\section{Data}
%\label{data}
%\section{Datasets}
\label{sec:arrhythmias:data}

\subsection*{Training}
We collect and annotate a dataset of 64,121 ECG records from 29,163 patients.
The ECG data is sampled at a frequency of 200 Hz and is collected from a
single-lead, noninvasive and  continuous monitoring device called the Zio Patch
which has a wear period up to 14 days \cite{turakhia2013diagnostic}. Each ECG
record in the training set is 30 seconds long and can contain more than one
rhythm type. Each record is annotated by a clinical ECG expert: the expert
highlights segments of the signal and marks it as corresponding to one of the
14 rhythm classes.

The 30 second records were annotated using a web-based ECG annotation tool
designed for this work. Label annotations were done by a group of Certified
Cardiographic Technicians who have completed extensive training in arrhythmia
detection and a cardiographic certification examination by Cardiovascular
Credentialing International. The technicians were guided through the interface
before they could annotate records. All rhythms present in a strip were labeled
from their corresponding onset to offset, resulting in full segmentation of the
input ECG data. To improve labeling consistency among different annotators,
specific rules were devised regarding each rhythm transition.

We split the dataset into a training and validation set. The training set
contains 90\% of the data. We split the dataset so that there is no patient
overlap between the training and validation sets (as well as the test set
described below).

\subsection*{Testing}

We collect a test set of 336 records from 328 unique patients. For the test
set, ground truth annotations for each record were obtained by a committee of
three board-certified cardiologists; there are three committees responsible for
different splits of the test set. The cardiologists discussed each individual
record as a group and came to a consensus labeling. For each record in the test
set we also collect 6 individual annotations from cardiologists not
participating in the group. This is used to assess performance of the model
compared to an individual cardiologist.

\subsection*{Rhythm Classes}
We identify 12 heart arrhythmias, sinus rhythm and noise for a total of 14
output classes. The arrhythmias are characterized by a variety of features.
Table~\ref{tab:rhythms} in the Appendix shows an example of each rhythm type we
classify. The noise label is assigned when the device is disconnected from the
skin or when the baseline noise in the ECG makes identification of the
underlying rhythm impossible.

The morphology of the ECG during a single heart-beat as well as the pattern of
the activity of the heart over time determine the underlying rhythm. In some
cases the distinction between the rhythms can be subtle yet critical for
treatment. For example two forms of second degree AV Block, Mobitz I
(Wenckebach) and Mobitz II (here referred to as AVB\_TYPE2) can be difficult to
distinguish. Wenckebach is considered benign and Mobitz II is considered
pathological, requiring immediate attention \cite{dubin1996rapid}. 

Table~\ref{tab:rhythms} in the Appendix also shows the number of unique
patients in the training (including validation) set and test set for each
rhythm type.

%
%\input{comparison}
%\section{Results}
%\label{results}
%\section{Results}
\label{sec:scaling_asr:results}

%
%\section{Analysis}
%\label{analysis}
%\section{Analysis}

The model outperforms the average cardiologist score on both the sequence and
the set F1 metrics. Figure~\ref{fig:arrhythmia:model_confusion} shows a
confusion matrix of the model predictions on the test set. Many arrhythmias are
confused with the sinus rhythm. We expect that part of this is due to the
sometimes ambiguous location of the exact onset and offset of the arrhythmia in
the ECG record.

\begin{figure}
\begin{subfigure}{.5\textwidth}
  \centering
  \includegraphics[width=0.9\linewidth]{arrhythmias/figures/model_confusions.pdf}
  \caption{Model}
  \label{fig:arrhythmia:model_confusion}
\end{subfigure}
\begin{subfigure}{.5\textwidth}
  \centering
  \includegraphics[width=0.9\linewidth]{arrhythmias/figures/human_confusions.pdf}
  \caption{Individual Cardiologists}
  \label{fig:arrhythmia:human_confusion}
\end{subfigure}
\caption{Confusion matrices for the model and individual cardiologist
         predictions with the committee consensus annotations as ground
         truth.}
\end{figure}

Often the mistakes made by the model are understandable. For example, confusing
Wenckebach and AVB\_Type2 makes sense given that the two rhythms in general
have very similar ECG morphologies. Similarly, Supraventricular Tachycardia
(SVT) and Atrial Fibrillation (AFIB) are often confused with Atrial Flutter
(AFL) which is understandable given that they are all atrial arrhythmias. We
also note that Idioventricular Rhythm (IVR) is sometimes mistaken as
Ventricular Tachycardia (VT), which again makes sense given that the two only
differ in heart-rate and are difficult to distinguish close to the 100 beats
per minute delineation.

One of the most common confusions is between Ectopic Atrial Rhythm (EAR) and
the sinus rhythm. The main distinguishing criteria for this rhythm is an
irregular P wave. This can be subtle to detect especially when the P wave has a
small amplitude or when noise is present in the signal.


%\begin{figure}[t]
%  \centering
%  \includegraphics[width=\linewidth]{figs/confusion.pdf}
%  \caption{
%    A confusion matrix for the model predictions on the test set. Many of the mistakes the model makes are not surprising. For example, confusing second degree AV Block (Type 2) with Wenckebach makes sense given the often similar expression of the two arrhythmias in the ECG record.
%  }
%  \label{fig:confusion}
%\end{figure}
%\section{Related Work}
%\label{related}
%\section{Related Work}
\label{sec:deepspeech:related}

Several parts of our work are inspired by previous results. Neural network
acoustic models and other connectionist approaches were first introduced to
speech pipelines in the early 1990s~\cite{bourlard93, renals1994, ellis1999}.
These systems, similar to DNN acoustic models~\cite{mohamed2011, hinton2012,
dahl2011a}, replace only one stage of the speech recognition pipeline.
Mechanically, our system is similar to other efforts to build end-to-end speech
systems from deep learning algorithms. For example,
Graves~et~al.~\cite{graves2006} have previously introduced the ``Connectionist
Temporal Classification'' (CTC) loss function for scoring transcriptions
produced by RNNs and, with LSTM networks, have previously applied this approach
to speech~\cite{graves2014}. We similarly adopt the CTC loss for part of our
training procedure but use much simpler recurrent networks with
rectified-linear activations~\cite{glorot2011, maas2013, nair2010}. Our
recurrent network is similar to the bidirectional RNN used by Hannun et
al.~\cite{hannun2014firstpass}, but with multiple changes to enhance its
scalability. By focusing on scalability, we have shown that these simpler
networks can be effective even without the more complex LSTM machinery.

Our work is certainly not the first to exploit scalability to improve
performance of DL algorithms. The value of scalability in deep learning is
well-studied~\cite{coates2011b, le2013} and the use of parallel processors
(including GPUs) has been instrumental to recent large-scale DL
results~\cite{szegedy2015, le2013}. Early ports of DL algorithms to GPUs
revealed significant speed gains~\cite{raina2009}. Researchers have also begun
choosing designs that map well to GPU hardware to gain even more efficiency,
including convolutional~\cite{krizhevsky2012imagenet, ciresan2011,
sainath2013deep} and locally connected~\cite{coates2013cotshpc, ciresan2012}
networks, especially when optimized libraries like
cuDNN~\cite{chetlur2014cudnn} and BLAS are available. Indeed, using
high-performance computing infrastructure, it is possible today to train neural
networks with more than 10 billion connections~\cite{coates2013cotshpc} using
clusters of GPUs. These results inspired us to focus first on making scalable
design choices to efficiently utilize many GPUs before trying to engineer the
algorithms and models themselves.

With the potential to train large models, there is a need for large training
sets as well. In other fields, such as computer vision, large labeled training
sets have enabled significant leaps in performance as they are used to feed
larger and larger DL systems~\cite{szegedy2105, krizhevsky2012imagenet}. In
speech recognition, however, such large training sets are less common, with
typical benchmarks having training sets ranging from tens of hours (e.g. the
Wall Street Journal corpus with 80 hours) to several hundreds of hours (e.g.
Switchboard and Broadcast News). Larger benchmark datasets, such as the Fisher
corpus~\cite{cieri2004fisher} with 2000 hours of transcribed speech, are rare
and only recently being studied. To fully utilize the expressive power of the
recurrent networks available to us, we rely not only on large sets of labeled
utterances, but also on synthesis techniques to generate novel examples. This
approach is well known in computer vision~\cite{sapp2008, lecun2004,
coates2011} but we have found this especially convenient and effective for
speech when done properly.


%\section{Conclusion}
%\label{conclusion}

%\section{Conclusion}

We presented a decoding algorithm which enables first-pass LVCSR with a
language model for CTC-trained neural networks. This decoding approach removes
the lingering dependence on HMM-based systems found in previous work.
Furthermore, first-pass decoding demonstrates the capabilities of a CTC-trained
system without the confounding factor of potential effects from pruning the
search space via a provided lattice.  While our results do not outperform the
best HMM-based systems on the WSJ corpus, they demonstrate the promise of
CTC-based speech recognition systems.

Our experiments with BRDNN further simplify the infrastructure needed to create
CTC-based speech recognition systems. The BRDNN is overall a less complex
architecture than LSTMs and can relatively easily be made to run on GPUs since
large matrix multiplications dominate the computation. However, our experiments
suggest that recurrent connections are critical for good performance.
Bi-directional recurrence helps beyond single direction recurrence but could be
sacrificed in cases that require low-latency, online speech recognition. Taken
together with previous work on CTC-based LVCSR, we believe there is an exciting
path forward for high quality LVCSR without the complexity of HMM-based
infrastructure.

%\input{rhythms}
